\documentclass[11pt,twocolumn]{article}
\usepackage[margin=1in]{geometry}
\usepackage{hyperref}
\usepackage{graphicx}
\usepackage{amsmath}
\usepackage{enumitem}
\usepackage{booktabs}
\usepackage{microtype}

\hypersetup{
  colorlinks=true,
  linkcolor=blue,
  urlcolor=blue,
  citecolor=blue
}

\title{\textbf{Eye-Tracking as a Game Input:\\
Porting Reflexio to the Web and Evaluating Gaze-Based Control}}

\author{Mengxiang Jiang\\
CSCE 624: Sketch Recognition
Spring 2026
}

\date{}

\begin{document}

\maketitle

% ============================================================
\section{Introduction and Motivation}
% ============================================================

Alternative human-computer input modalities have long fascinated researchers and
hobbyists alike. Moving beyond the conventional keyboard-and-mouse or gamepad
paradigm opens the door to more accessible, expressive, or simply novel ways to
interact with software. My interest in non-traditional game input began during
my undergraduate studies, where I built a system for playing Pong using
electroencephalography (EEG) signals recorded from a consumer headset
\cite{eeg_pong}. That project demonstrated that exotic input channels can be
made functional, but also highlighted the significant gap between
``technically works'' and ``pleasant to use.''

This project revisits that question in a different context: gaze-based control
for a puzzle platformer. The game chosen is \textit{Reflexio}, a tile-based
puzzle platformer I originally wrote as an undergraduate course project in C\#
using the XNA/FNA framework \cite{reflexio_xna}. Reflexio's signature mechanic
is \textit{reflection}: the player can mirror the entire level, or a subset of
rows, columns, or diagonal bands, along a selectable line. Because choosing
\emph{which} reflection line to activate is a discrete, deliberate selection
task rather than a continuous pointing task, it is a plausible candidate for
gaze control; the player looks at the indicator for the desired axis and then
uses a separate button press to trigger the reflection.

The goals of this project are therefore threefold: (1) port Reflexio to a
web-accessible platform so that participants can play without installing
anything; (2) integrate a webcam-based eye-tracking input channel for the
reflection-line selection sub-task; and (3) collect qualitative feedback on
whether the gaze control improves, impairs, or simply does not affect the
playing experience.

% ============================================================
\section{Related Background}
% ============================================================

\subsection{Eye-Tracking in HCI and Games}

Eye-tracking has been studied as a computer input modality since at least the
1990s \cite{jacob1991}. The core challenge in using gaze as a \emph{command}
channel is the \textbf{Midas touch problem} \cite{jacob1990}: because users
look at many things they do not intend to select, a purely gaze-triggered
system fires spuriously. The canonical mitigation is to require a secondary
confirmation gesture (a button press, dwell timer, or blink) before any action
is executed. This project adopts the button-press approach: the player uses
gaze to navigate among reflection-line indicators, and a separate key or button
to execute the reflection.

\subsection{Saccades and Smooth Pursuit}

Eye movements fall into two primary categories relevant here. \textbf{Saccades}
are rapid, ballistic jumps between fixation points, typically lasting 20--200
ms. \textbf{Fixations} are the relatively stable periods between saccades,
during which the visual system processes information. Consumer webcam-based
trackers such as WebGazer.js operate at video frame rates (typically 30 fps)
and return the estimated screen coordinate of the current fixation. The
relatively coarse temporal resolution means that saccade onset and offset cannot
be detected precisely; instead, a change in fixation region must be inferred
from a sustained shift in the gaze estimate.

\subsection{Fitts's Law}

Fitts's law \cite{fitts1954} predicts the time required to move a pointer to a
target as a function of the distance to the target and its width:
$T = a + b \log_2\!\left(1 + D/W\right)$. For gaze-based pointing, the
effective ``width'' of a target is limited by tracker accuracy rather than
target size, which severely penalises small or closely spaced targets. The
reflection-line indicators in Reflexio are relatively small, making gaze
pointing challenging even in principle.

\subsection{Kalman Filtering for Gaze Smoothing}

Raw webcam-based gaze estimates are noisy. A scalar Kalman filter
\cite{kalman1960} is a standard tool for smoothing a one-dimensional noisy
measurement sequence under a linear-Gaussian model. Applying independent
filters to the $x$ and $y$ gaze coordinates reduces high-frequency jitter
while introducing a small lag proportional to the process-noise-to-measurement-noise
ratio. This project applies a 1-D Kalman filter to each axis of the WebGazer
output.

\subsection{WebGazer.js}

WebGazer.js \cite{webgazer} is an open-source JavaScript library developed at
Brown University that performs real-time gaze estimation using only a standard
webcam accessed through the browser's \texttt{getUserMedia} API. It uses facial
landmark detection to locate the eyes and trains a regression model (ridge
regression by default) mapping eye-region image features to screen coordinates.
Training data is collected implicitly as the user interacts with the page
(click and cursor positions serve as ground-truth labels). No specialised
hardware is required, making it well suited for a web-deployed experiment.

% ============================================================
\section{Method and Approach}
% ============================================================

The project proceeded in two phases.

\textbf{Phase 1 -- Web port.} The original Reflexio codebase targets the
XNA/FNA desktop runtime and cannot run in a browser. The first task was
therefore to produce a faithful web port. After evaluating several JavaScript
game frameworks, Phaser 3 \cite{phaser} was selected for its built-in Matter.js
physics engine, rich sprite and animation support, and active community. The
port preserves all 53 published levels, the core reflection mechanic, player
movement and physics, key/door/switch puzzles, and the main menu and level
selection flow. Mobile touch controls were added later in response to user
feedback (Section~\ref{sec:results}).

\textbf{Phase 2 -- Gaze input for reflection selection.} The reflection-line
selection sub-task was identified as the most suitable target for gaze input
because it is a discrete choice among a small number of options (horizontal,
vertical, or diagonal axes) and the Midas touch problem is naturally mitigated:
looking at an indicator changes the \emph{highlighted} line but does not
execute the reflection; a separate button press is required. This matches the
intent-plus-confirmation pattern recommended in the gaze-interaction literature.

Gaze is used to cycle through reflection line categories via detected saccades.
A large horizontal saccade (displacement exceeding a calibrated threshold)
cycles the horizontal reflection lines; a large vertical saccade cycles the
vertical lines. Diagonal lines are selected by keyboard or gamepad as a
fallback. An auto-cycling feature was added so that holding a saccade direction
repeatedly advances through the available lines without requiring repeated
saccades.

% ============================================================
\section{Implementation Details}
% ============================================================

\subsection{Game Port}

The port was implemented in JavaScript using Phaser 3 and bundled with Vite.
The original C\# codebase used reflection (in the language sense) to
instantiate game objects from XML level files; this was replicated in JavaScript
via a class registry that maps XML tag names to constructors. All 64 level XML
files from the original ship unchanged. Physics bodies are provided by
Matter.js: walls and most objects use static rectangle bodies, the player uses
a 20-vertex polygon approximating an ellipse (matching the original's
\texttt{PolygonTools.CreateEllipse}), and grabbable blocks use dynamic bodies
connected to the player via a distance constraint.

A significant portion of the implementation work involved learning JavaScript
and the Phaser API. Claude (Anthropic's AI assistant) was used extensively as a
coding aid throughout the port, providing API guidance, debugging support, and
code generation for boilerplate-heavy sections.

\subsection{Touch and Mobile Controls}

Following a request from a friend who discovered the game on his phone, a DOM
overlay providing virtual controls was added. The overlay uses Pointer Events
(which unify mouse and touch input) and positions all elements with CSS
\texttt{calc()} expressions in viewport units so that the layout adapts to any
screen size without JavaScript measurement. Two D-pads handle movement and
reflection-line navigation; A and B buttons handle jump and reflect. The
browser's Fullscreen API and Screen Orientation API are invoked on the first
user gesture to maximise the playable area and lock landscape orientation on
mobile devices.

\subsection{Eye-Tracking Integration}

WebGazer.js is loaded on demand from a CDN when the user enables the eye
tracker. After \texttt{webgazer.begin()} acquires camera permission and starts
the regression model, raw gaze coordinates are post-processed by two
independent 1-D Kalman filters (one per axis). The filter parameters were
tuned empirically to balance lag against noise reduction.

The smoothed gaze point is compared each frame to a running baseline (an
exponential moving average of recent gaze positions). When the displacement
from the baseline exceeds a threshold---100 pixels horizontally or 60 pixels
vertically---a saccade is registered and the corresponding reflection-line
category is advanced. A 120-frame cooldown prevents multiple triggers from a
single saccade. Outlier rejection discards any single-frame jump exceeding 300
pixels, which typically corresponds to tracker loss or a blink rather than a
genuine saccade.

Blink detection is implemented by counting consecutive null callbacks from
WebGazer (indicating the eye region is occluded). A run of 2--18 null samples
is classified as a blink. Blinks were initially considered as a potential
command channel but were ultimately not used due to reliability concerns.

The webcam preview and WebGazer's built-in prediction-point overlay are
displayed in a fixed corner of the screen at reduced opacity so the player can
monitor tracker state without it dominating the view. Eye-tracking buttons
(toggle and calibrate) are suppressed on mobile browsers, where webcam-based
gaze tracking is impractical, using a user-agent check.

\subsection{Kalman Filter}

The 1-D Kalman filter implementation maintains state $(\hat{x}, P)$ representing
the estimated gaze coordinate and its error covariance. Each frame the
prediction step advances the state with process noise $Q = 0.1$ and the
update step incorporates the WebGazer measurement with measurement noise
$R = 30$. These values were chosen so that the filter follows genuine fixation
shifts within roughly five frames while attenuating single-frame noise spikes.

% ============================================================
\section{Results and Evaluation}
\label{sec:results}
% ============================================================

\subsection{Tracking Accuracy}

Initial testing of the raw WebGazer output revealed gaze estimation errors
exceeding 100 pixels on a 1920$\times$1080 display. The reflection-line
indicators in Reflexio are small (on the order of 30--40 pixels), making
accurate pixel-level selection infeasible. Applying the Kalman filter reduced
the effective jitter to approximately 50 pixels, which is an improvement but
still insufficient for reliable indicator selection by fixation alone.

\subsection{Saccade-Based Selection}

Switching from fixation-based to saccade-based selection significantly reduced
the accuracy demands: the user need only produce a large-magnitude eye movement
in the correct cardinal direction, rather than fixating a specific small target.
In informal testing this was more reliable than fixation-based selection.
However, the discretisation (horizontal saccade $\to$ cycle horizontal lines;
vertical saccade $\to$ cycle vertical lines) means the user must repeat the
same saccade multiple times to step through several lines of the same
orientation, which is fatiguing and inconsistent. The auto-cycling feature
(advancing automatically every two seconds after a saccade is registered)
partially mitigates this but introduces unwanted advances when the player is not
actively selecting.

\subsection{User Feedback}

The author and three friends played Reflexio with eye-tracking enabled. The
unanimous reaction was that the gaze control made the game noticeably harder and
less enjoyable to play. Specific complaints included:

\begin{itemize}[noitemsep]
  \item Calibration drift: tracking quality degraded over a session,
        requiring frequent manual recalibration.
  \item Cognitive load: monitoring gaze direction in addition to the in-game
        state was mentally taxing.
  \item Misdirected saccades: glancing at score displays or the webcam
        preview triggered unintended line changes.
  \item Inconsistency: the same intentional saccade did not always produce the
        same result.
\end{itemize}

One participant who tested the game on a mobile phone noted the absence of touch
controls. This feedback directly motivated the implementation of the virtual
D-pad overlay described in Section~4.2, which was well received.

% ============================================================
\section{Discussion and Limitations}
% ============================================================

The central finding is that eye-tracking input, at least as implemented here
with consumer-grade webcam-based tracking, is not a viable replacement for
keyboard or gamepad input in a fast-paced puzzle platformer. The accuracy
limitations of WebGazer.js on typical hardware are a fundamental constraint
that no amount of software-side smoothing can fully overcome; the technique
requires purpose-built hardware (infrared-illuminated eye trackers with
sub-pixel accuracy) to reach the precision needed for reliably selecting small
on-screen targets.

The choice of Reflexio's reflection mechanic as the target sub-task was
deliberate and reasonable: it is a discrete choice task, it naturally avoids
the Midas touch problem, and it does not require continuous high-bandwidth
gaze control. Nevertheless, even this limited use of gaze proved frustrating in
practice, primarily due to calibration drift and the difficulty of producing
consistent, repeatable saccades.

Several limitations constrain the generalisability of these results. The
evaluation was informal, with three participants, no control condition, and no
quantitative performance metrics such as level completion time or error rate.
A more rigorous study would include a within-subjects comparison of keyboard-only
versus keyboard-plus-gaze conditions, objective performance measures, and a
larger and more diverse participant pool. The port itself, while faithful in
mechanics, has not been systematically validated against the original for
physics fidelity across all 53 levels.

% ============================================================
\section{Conclusion}
% ============================================================

This project produced a complete web port of the puzzle platformer Reflexio,
playable at \url{https://kumj2028.github.io/reflexio-phaser/} with full mobile
touch control support, and integrated webcam-based eye tracking for the
reflection-line selection sub-task. The eye-tracking integration demonstrated
that novelty of input modality does not imply usability: participants
consistently found the gaze control more burdensome than the standard keyboard
alternative. The technical limitations of consumer-grade webcam tracking---high
noise, calibration drift, and sensitivity to ambient lighting---impose accuracy
constraints that are difficult to overcome in software. The project reinforces
a well-established principle in HCI: a new control scheme must improve the user
experience in a measurable way; if the existing scheme works well, replacing it
with a novel but inferior alternative degrades rather than enhances the game.
The most practically useful outcome of the project was not the eye-tracking
integration itself but the mobile touch control overlay, which extended the
game's accessibility to smartphone users and received positive feedback.

Source code is available at \url{https://github.com/kumj2028/reflexio-phaser}.

% ============================================================
\begin{thebibliography}{9}

\bibitem{eeg_pong}
C.\ Moyes and M.\ Jiang,
``EEG-Controlled Pong,''
ECE 4760 Final Project, Cornell University, Spring 2012.
\url{https://people.ece.cornell.edu/land/courses/ece4760/FinalProjects/s2012/cwm55/cwm55_mj294/}

\bibitem{reflexio_xna}
M.\ Jiang,
``Reflexio,''
CS 3152 course project, Cornell University.
\url{https://kumj2028.github.io/CS3152/index.html}

\bibitem{jacob1991}
R.\ J.\ K.\ Jacob,
``The use of eye movements in human-computer interaction techniques: What you look at is what you get,''
\textit{ACM Transactions on Information Systems}, vol.~9, no.~2, pp.~152--169, 1991.

\bibitem{jacob1990}
R.\ J.\ K.\ Jacob,
``What you look at is what you get: eye movement-based interaction techniques,''
in \textit{Proc.\ ACM CHI}, 1990, pp.~11--18.

\bibitem{fitts1954}
P.\ M.\ Fitts,
``The information capacity of the human motor system in controlling the amplitude of movement,''
\textit{Journal of Experimental Psychology}, vol.~47, no.~6, pp.~381--391, 1954.

\bibitem{kalman1960}
R.\ E.\ Kalman,
``A new approach to linear filtering and prediction problems,''
\textit{Journal of Basic Engineering}, vol.~82, no.~1, pp.~35--45, 1960.

\bibitem{webgazer}
A.\ Papoutsaki, P.\ Sangkloy, J.\ Laskey, N.\ Daskalova, J.\ Huang, and J.\ Hays,
``WebGazer: Scalable webcam eye tracking using user interactions,''
in \textit{Proc.\ IJCAI}, 2016.
\url{https://webgazer.cs.brown.edu/}

\bibitem{phaser}
Photon Storm Ltd.,
``Phaser 3,'' 2024.
\url{https://phaser.io/}

\end{thebibliography}

\end{document}
